% Figure: diffraction limit versus atmospheric seeing as a function of aperture.
% The diffraction-limited resolution 1.22 lambda/D falls as 1/D (it improves with
% aperture), while ground-level atmospheric seeing is a fixed floor near 1 arcsec
% independent of D. The two curves cross near D = 0.14 m at 550 nm; beyond that
% the ground telescope is seeing-limited, not diffraction-limited, until adaptive
% optics restores the diffraction curve.
\begin{figure}[htbp]
\centering
\begin{tikzpicture}
\begin{axis}[
  width=13cm, height=7.2cm,
  xlabel={primary aperture $D$ (m)},
  ylabel={angular resolution (arc-second)},
  xmode=log, ymode=log,
  xmin=0.05, xmax=40, ymin=0.003, ymax=4,
  axis lines=left,
  xtick={0.1,0.2,0.5,1,2,5,10,20},
  xticklabels={0.1,0.2,0.5,1,2,5,10,20},
  ytick={0.005,0.01,0.05,0.1,0.5,1,2},
  yticklabels={0.005,0.01,0.05,0.1,0.5,1,2},
  tick label style={font=\scriptsize},
  label style={font=\small},
  legend style={font=\scriptsize, at={(0.98,0.97)}, anchor=north east, draw=none},
  samples=80, domain=0.05:40,
]
% diffraction limit in arcsec: theta = 1.22 * 550e-9 / D * (180/pi)*3600
% = 1.22*550e-9*206265/D = 0.1385 / D
\addplot[engblue, very thick] {0.1385/x};
\addlegendentry{diffraction limit $1.22\lambda/D$}
% atmospheric seeing floor ~ 1 arcsec
\addplot[engred, thick, dashed] {1.0};
\addlegendentry{typical seeing ($\sim 1''$)}
% adaptive-optics restored region marker (vertical guide at 8 m)
\draw[enggray, dotted] (axis cs:8,0.003) -- (axis cs:8,4);
\node[enggray, font=\scriptsize, anchor=south, rotate=90] at (axis cs:8,0.02)
  {$8\,$m class};
% crossover marker
\node[engblack, font=\scriptsize, anchor=west] at (axis cs:0.16,1.25)
  {crossover $\approx 0.14\,$m};
\draw[engblack, ->] (axis cs:0.16,1.1) -- (axis cs:0.14,1.0);
\node[engorange, font=\scriptsize, anchor=south west] at (axis cs:1.1,0.013)
  {AO target with correction};
\end{axis}
\end{tikzpicture}
\caption[Diffraction versus seeing]{The diffraction-limited resolution
$1.22\lambda/D$ improves as $1/D$ with aperture (sloping line), while
ground-level atmospheric seeing is a fixed floor near one arc-second
independent of aperture (dashed line). The two cross near $D\approx
0.14\,\si{\meter}$ at $550\,\si{\nano\meter}$: a backyard refractor smaller
than that is diffraction-limited, while every larger telescope is
seeing-limited at a typical site and delivers a one-arc-second image whatever
its aperture. Adaptive optics measures and cancels the atmospheric phase error
in real time and pulls the delivered resolution back down toward the blue
diffraction curve, which is why an $8\,\si{\meter}$ telescope is worth building
even though the raw atmosphere would waste its aperture.}
\label{fig:vol04ch11:seeing-vs-diffraction}
\end{figure}
